% SGA 3，报告 IX，中文重构。
% 范围：第 2、3 节全文。
% 权威底本：Polo--Gille Exp9-8nov09.pdf，组件本地第 4--10 页
% （完整阅读本第 650--656 页），止于第 4 节之前。
% Jacob Reinhold 的 CC BY 4.0 Markdown 谱系仅用于措辞比对。
% 文字、公式、编号、注释和引用均以 Polo--Gille PDF 为准。

\setcounter{footnote}{2}

% 源页：Exp9 组件本地第 4 页；完整阅读本第 650 页；原版印刷第 41--42 页。
\SGARefTarget{sga3:ix:section:2}\section*{2. 将可对角化群的某些性质推广到乘法型群}
\addcontentsline{toc}{section}{2. 将可对角化群的某些性质推广到乘法型群}

鉴于\SGARefLink{sga3:ix:definition:1:1}{1.1 中的定义}以及所考虑结果的“局部”性质（对忠实平坦拟紧拓扑而言），这些推广本质上只是前一篇报告结果的直接推论。

\medskip
\noindent\SGARefTarget{sga3:ix:definition:2:0}\textbf{定义 2.0.}
\footnote{编者注：此定义为新增；它出现在\SGARefLink{sga3:ix:proposition:2:3}{命题~2.3}和\SGARefLink{sga3:ix:proposition:2:7}{2.7}中。}
\textit{用 \(\mathcal M\) 表示忠实平坦拟紧态射所成的集合。}

\medskip
\noindent\SGARefTarget{sga3:ix:proposition:2:1}\textbf{命题 2.1.}
\textit{设 \(G\) 是预概形 \(S\) 上的乘法型群。则：}
\begin{enumerate}
  \item[(a)]\SGARefTarget{sga3:ix:proposition:2:1:item:a:01}
  \textit{\(G\) 在 \(S\) 上忠实平坦，且在 \(S\) 上仿射（因而在 \(S\) 上拟紧）；}

  \item[(b)]\SGARefTarget{sga3:ix:proposition:2:1:item:b:01}
  \textit{\(G\) 在 \(S\) 上有限型，当且仅当 \(G\) 在 \(S\) 上有限呈示，也当且仅当对每个 \(s\in S\)，\(G\) 在 \(s\) 处的类型由一个有限型交换群给出；}

  \item[(c)]\SGARefTarget{sga3:ix:proposition:2:1:item:c:01}
  \textit{\(G\) 在 \(S\) 上有限，当且仅当对每个 \(s\in S\)，\(G\) 在 \(s\) 处的类型由有限交换群给出，也当且仅当（在 \(S\) 拟紧时）\(G\) 在 \(S\) 上有限型，且被某个整数 \(n>0\) 零化；}

  \item[\((c')\)]\SGARefTarget{sga3:ix:proposition:2:1:item:cprime:01}
  \textit{\(G\) 在 \(S\) 上整，当且仅当对每个 \(s\in S\)，\(G\) 在 \(s\) 处的类型由挠交换群给出；}

  \item[(d)]\SGARefTarget{sga3:ix:proposition:2:1:item:d:01}
  \textit{\(G\) 是单位群，当且仅当对每个 \(s\in S\)，\(G\) 在 \(s\) 处的类型由零交换群给出；}

  \item[(e)]\SGARefTarget{sga3:ix:proposition:2:1:item:e:01}
  \textit{\(G\) 在 \(S\) 上平滑，当且仅当对每个 \(s\in S\)，\(G\) 在 \(s\) 处的类型由有限型交换群给出，且该交换群的挠子群的阶与 \(\kappa(s)\) 的特征互素。}
\end{enumerate}

这些断言由 \SGARefLink{sga3:viii:proposition:2:1}{VIII~2.1} 推出，这里须注意有关性质可沿忠实平坦拟紧态射下降（参见 SGA~1，VIII，或 EGA~IV\textsubscript{2}，\S2）。

利用 \SGARefLink{sga3:viii:corollary:3:5}{VIII~3.5}，类似地得到：

\medskip
\noindent\SGARefTarget{sga3:ix:proposition:2:2}\textbf{命题 2.2.}
\textit{设 \(G\) 是 \(S\) 上的有限型乘法型群。对每个整数 \(n\ne0\)，核 \({}_nG=\ker(n\cdot\operatorname{id}_G)\) 是在 \(S\) 上有限的乘法型群。}

\medskip
\noindent\SGARefTarget{sga3:ix:proposition:2:3}\textbf{命题 2.3.}
\textit{设 \(G\) 是预概形 \(S\) 上的乘法型群，自由地右作用于一个在 \(S\) 上仿射的 \(S\)-预概形 \(X\)。则：}
\begin{enumerate}
  \item[(a)]\SGARefTarget{sga3:ix:proposition:2:3:item:a:01}
  \textit{\(G\) 在 \(X\) 中定义的等价关系是 \(\mathcal M\)-有效的（\SGARefLink{sga3:iv:subsection:3:4}{IV~3.4}），且 \(Y=X/G\) 在 \(S\) 上仿射；}

  \item[(b)]\SGARefTarget{sga3:ix:proposition:2:3:item:b:01}
  \textit{此外，若 \(X\) 在 \(S\) 上有限呈示（分别地，有限型），则 \(Y\) 也如此。}
\end{enumerate}

\SGARefLink{sga3:ix:proposition:2:3:item:a:01}{第一个断言} 由处理 \(G\) 可对角化情形的 \SGARefLink{sga3:viii:theorem:5:1}{VIII~5.1}，以及允许归结到该情形的 \SGARefLink{sga3:iv:proposition:3:5:2}{IV~3.5.2} 推出。事实上，忠实平坦拟紧态射 \(S'\to S\) 是概形上的仿射概形所成纤维范畴的有效下降态射：若 \(Y'\) 在 \(S'\) 上仿射，并赋有相对于 \(S'\to S\) 的下降数据，则此数据有效，即 \(Y'\) 来自一个在 \(S\) 上仿射的概形 \(Y\)（参见 SGA~1，\SGARefLink{sga3:viii:proposition:2:1}{VIII~2.1}）。

% 源页：Exp9 组件本地第 5 页；完整阅读本第 651 页；原版印刷第 43 页。
对于\SGARefLink{sga3:ix:proposition:2:3:item:b:01}{第二个断言}，同样可归结到 \SGARefLink{sga3:viii:corollary:5:8}{VIII~5.8} 的可对角化情形，因为所考虑的有限性条件可沿忠实平坦拟紧态射下降（SGA~1，VIII~3.3 和 3.6\footnote{编者注：另见 V，9.1，或 EGA~IV\textsubscript{2}，2.7.1。}）。仿照 VIII 的推论~5.5--5.7，可由 \SGARefLink{sga3:ix:proposition:2:3}{2.3} 推出：

\medskip
\noindent\SGARefTarget{sga3:ix:corollary:2:4}\textbf{推论 2.4.}
\textit{在 \SGARefLink{sga3:ix:proposition:2:3}{2.3} 的假设下，图态射}
\SGARefTarget{sga3:ix:corollary:2:4:diagram:01}
\[
  X\times_S G\longrightarrow X\times_S X
\]
\textit{是闭浸入。对 \(X\) 在 \(S\) 上的每个截面 \(\sigma\)，对应的态射}
\[
  G\longrightarrow X,\qquad g\longmapsto\sigma\cdot g,
\]
\textit{是闭浸入。}

特别地：

\medskip
\noindent\SGARefTarget{sga3:ix:corollary:2:5}\textbf{推论 2.5.}
\textit{设 \(u:G\to H\) 是群 \(S\)-预概形的单态射，其中 \(G\) 是乘法型，而 \(H\) 在 \(S\) 上仿射。则 \(u\) 是闭浸入，\(H/G=Y\) 存在且在 \(S\) 上仿射，\(H\) 是 \(Y\) 上在 \(G_Y\) 作用下的主齐性丛。}

\medskip
\noindent\SGARefTarget{sga3:ix:remark:2:6}\textbf{注记 2.6.}
设 \(u:G\to H\) 是群 \(S\)-预概形的单态射，其中 \(G\) 是 \(S\) 上的有限型乘法型群，而 \(H\) 在 \(S\) 上有限呈示且分离。此时利用 \SGARefLink{sga3:viii:proposition:7:12}{VIII~7.12} 可证明 \(u\) 是闭浸入（见注记~\SGARefLink{sga3:viii:remark:7-13:item:b}{VIII~7.13(b)}）。

\medskip
\noindent\SGARefTarget{sga3:ix:proposition:2:7}\textbf{命题 2.7.}
\textit{设 \(u:G\to H\) 是乘法型 \(S\)-群的同态，且 \(H\) 在 \(S\) 上有限型。则：}
\begin{enumerate}
  \item[(i)]\SGARefTarget{sga3:ix:proposition:2:7:item:i:01}
  \textit{\(G_0=\ker u\) 是乘法型 \(S\)-群；\(G_0\) 在 \(G\) 中定义的等价关系是 \(\mathcal M\)-有效的，因此（\SGARefLink{sga3:iv:subsection:3:4}{IV~3.4}）\(u\) 分解为}
  \SGARefTarget{sga3:ix:proposition:2:7:diagram:01}
  \[
    G\longrightarrow G/G_0=I\longrightarrow H,
  \]
  \textit{其中 \(I\to H\) 是 \(S\)-群的单态射\footnote{编者注：事实上，甚至是闭浸入。}。群 \(I\) 是乘法型，\(I\) 在 \(H\) 中定义的等价关系是 \(\mathcal M\)-有效的，故 \(H_0=H/I\) 存在且也为乘法型；}

  \item[(ii)]\SGARefTarget{sga3:ix:proposition:2:7:item:ii:01}
  \textit{\(H_0\) 和 \(I\) 有限型，而当 \(G\) 有限型时，\(G_0\) 也有限型。}
\end{enumerate}

\emph{证明。}
仿照 \SGARefLink{sga3:ix:proposition:2:3}{2.3}，可归结到 \(G\) 与 \(H\) 可对角化的情形；这时 \SGARefLink{sga3:ix:proposition:2:7}{2.7} 归结为 \SGARefLink{sga3:viii:corollary:3:4}{VIII~3.4}。

记下以下推论。

% 源页：Exp9 组件本地第 6 页；完整阅读本第 652 页；原版印刷第 44 页。
\medskip
\noindent\SGARefTarget{sga3:ix:corollary:2:8}\textbf{推论 2.8.}
\begin{enumerate}
  \item[(a)]\SGARefTarget{sga3:ix:corollary:2:8:item:a:01}
  \textit{设 \(S\) 是预概形。有限型乘法型 \(S\)-群的范畴是阿贝尔范畴。}

  \item[(b)]\SGARefTarget{sga3:ix:corollary:2:8:item:b:01}
  \textit{设 \(u:G\to H\) 是此范畴中的同态。它是此范畴中的单态射（分别地，满态射，分别地，同构），当且仅当 \(u\) 是预概形的单态射（分别地，\(u\) 忠实平坦，分别地，\(u\) 是预概形的同构）。}
\end{enumerate}

只须注意，两个乘法型 \(S\)-群之积仍为乘法型（这是直接的），而且若两个因子均在 \(S\) 上有限型，其积显然也如此。\SGARefLink{sga3:ix:corollary:2:8}{2.8} 的其余部分由 \SGARefLink{sga3:ix:proposition:2:7}{2.7} 立即得到；细节留给读者。

\medskip
\noindent\SGARefTarget{sga3:ix:corollary:2:9}\textbf{推论 2.9.}
\textit{设 \(u:G\to H\) 是有限型乘法型 \(S\)-群的同态。设 \(U\) 为使 \(u_s:G_s\to H_s\) 是单态射（分别地，忠实平坦，分别地，同构）的点 \(s\in S\) 所成的集合。则 \(U\) 既开又闭，且诱导的同态 \(u_U:G_U\to H_U\) 是单态射（分别地，忠实平坦，分别地，同构）。}

设 \(P\)（分别地，\(Q\)）为 \(u\) 的核（分别地，余核）。由 \SGARefLink{sga3:ix:proposition:2:7}{2.7}，\(Q\) 存在，且 \(P\) 和 \(Q\) 是乘法型；此外，\(P\) 和 \(Q\) 的形成与每个基变换 \(S'\to S\) 交换，特别地，与取纤维交换。另一方面，\(u\) 是单态射（分别地，忠实平坦，分别地，同构），当且仅当 \(P=0\)（分别地，\(Q=0\)，分别地，\(P=Q=0\)）。这些观察将问题归结为验证以下断言：若 \(R\) 是乘法型 \(S\)-群，则使 \(R_s=0\) 的点 \(s\in S\) 所成的集合 \(U\) 既开又闭，且 \(R_U=0\)。这包含在 \SGARefLink{sga3:ix:remark:1:4:1:item:a}{注记~1.4.1(a)} 中。

\medskip
\noindent\SGARefTarget{sga3:ix:corollary:2:10}\textbf{推论 2.10.}
\textit{设 \(G\) 是有限型乘法型 \(S\)-群，\(H\) 和 \(H'\) 是两个乘法型子群。}
\begin{enumerate}
  \item[(a)]\SGARefTarget{sga3:ix:corollary:2:10:item:a:01}
  \textit{\(H''=H\cap H'\) 是 \(G\) 的乘法型子群。}

  \item[(b)]\SGARefTarget{sga3:ix:corollary:2:10:item:b:01}
  \textit{设 \(U\) 为满足 \(H_s\subseteq H'_s\)（分别地，\(H_s=H'_s\)）的 \(s\in S\) 所成的集合。则 \(U\) 既开又闭，且 \(H_U\subseteq H'_U\)（分别地，\(H_U=H'_U\)）。}
\end{enumerate}

当然，\(H\cap H'\) 中的交号表示函子意义下的交，即 \(H\times_GH'\)，它显然是 \(G\) 的 \(S\)-子群。同样，包含号（分别地，等号）表示 \(H\) 的子函子之间的序关系（分别地，相等），并非底层集合的包含（分别地，相等）。

先对包含态射 \(H\to G\) 应用 \SGARefLink{sga3:ix:proposition:2:7}{2.7}，可知 \(H\) 有限型；\(H'\) 也如此。于是由 \SGARefLink{sga3:ix:corollary:2:8}{2.8}，\(H\cap H'\) 是有限型乘法型群。（注意，从 \SGARefLink{sga3:ix:corollary:2:8}{2.8} 所考虑范畴到范畴 \((\mathbf{Sch})/S\) 的典范函子与有限逆极限交换。）

\(H\cap H'\) 的形成与基扩张交换，特别地，与取纤维交换。此外，\(H\subseteq H'\)（分别地，\(H=H'\)）等价于 \(H''=H\)（分别地，\(H''=H\) 和 \(H''=H'\)）。考虑典范同态 \(H''\to H\) 和 \(H''\to H'\)。则 \(U\) 是使诱导同态 \(H''_s\to H_s\) 为同构（分别地，使 \(H''_s\to H_s\) 和 \(H''_s\to H'_s\) 为同构）的点 \(s\in S\) 所成的集合。

% 源页：Exp9 组件本地第 7 页；完整阅读本第 653 页；原版印刷第 45--46 页。
由 \SGARefLink{sga3:ix:corollary:2:9}{2.9}，得到 \(U\) 既开又闭，且 \(H''_U\to H_U\)（分别地，\(H''_U\to H_U\) 和 \(H''_U\to H'_U\)）是同构。所要的结论由此得到。

\medskip
\noindent\SGARefTarget{sga3:ix:proposition:2:11}\textbf{命题 2.11.}
\textit{设 \(S\) 是预概形，\(G\) 是在 \(S\) 上有限型的乘法型 \(S\)-群，\(H\) 是 \(G\) 的乘法型子群，并置 \(K=G/H\)，它是 \(G\) 的乘法型商群。}
\begin{enumerate}
  \item[(i)]\SGARefTarget{sga3:ix:proposition:2:11:item:i:01}
  \textit{假设 \(G\) 平凡，即可对角化。存在 \(S\) 为既开又闭子集 \(S_i\) 的分划，使对每个 \(i\)，\(H_{S_i}\) 和 \(K_{S_i}\) 均可对角化。特别地，若 \(S\) 连通，则 \(H\) 和 \(K\) 可对角化。}

  \item[(ii)]\SGARefTarget{sga3:ix:proposition:2:11:item:ii:01}
  \textit{将“可对角化”换成“等平凡”，\SGARefLink{sga3:ix:proposition:2:11:item:i:01}{相同断言} 仍成立，只要 \(S\) 连通、局部连通或拟紧。}

  \item[(iii)]\SGARefTarget{sga3:ix:proposition:2:11:item:iii:01}
  \textit{若 \(G\) 局部平凡（分别地，局部等平凡，分别地，拟等平凡），则 \(K\) 和 \(H\) 也如此。}
\end{enumerate}

\emph{证明。}
\SGARefLink{sga3:ix:proposition:2:11:item:i:01}{(i)} 由假设，\(G=D_S(M)\)，其中 \(M\) 是有限型交换群。对 \(M\) 的每个商群 \(M_i\)，设 \(H_i=D_S(M_i)\) 为 \(G\) 的相应可对角化子群（\SGARefLink{sga3:viii:theorem:3:1}{VIII~3.1}）。设 \(S_i\) 为使 \(H_s=(H_i)_s\) 的点 \(s\in S\) 所成的集合。由 \SGARefLink{sga3:ix:corollary:2:10}{2.10}，\(S_i\) 既开又闭，且 \(H_{S_i}=(H_i)_{S_i}\)；因此 \(H_{S_i}\) 可对角化，从而 \(K_{S_i}\) 也可对角化（\SGARefLink{sga3:viii:theorem:3:1}{VIII~3.1}）。各 \(S_i\) 显然两两不交，我们断言它们覆盖 \(S\)。事实上，对每个 \(s\)，群 \(H_s\) 作为可对角化群 \(G_s\) 的子群，是可对角化的（参见下文 \SGARefLink{sga3:ix:proposition:8:1}{8.1}\footnote{编者注：已修正错误的引用 IV~4.7.5。注意，本报告的\SGARefLink{sga3:ix:section:8}{第~8~节}不依赖第~\SGARefLink{sga3:ix:section:3}{3}--\SGARefLink{sga3:ix:section:7}{7}~节。}）。因此由 \SGARefLink{sga3:viii:corollary:1:5}{VIII~1.5} 和 \SGARefLink{sga3:ix:theorem:3:2:item:b}{3.2(b)}，\(H_s\) 形如 \(D_{\kappa(s)}(M_i)\)。限于非空 \(S_i\) 所成的族，就得到断言 \SGARefLink{sga3:ix:proposition:2:11:item:i:01}{(i)}。

\SGARefLink{sga3:ix:proposition:2:11:item:ii:01}{(ii)} 由假设，存在有限平展满射 \(S'\to S\)，使 \(G_{S'}\) 可对角化。因此 \(S'\) 的每个点都有既开又闭邻域 \(U'\)，使 \(H_{U'}\) 和 \(K_{U'}\) 可对角化。\(U'\) 在 \(S\) 中的像 \(U\) 既开又闭，且 \(S'\to S\) 仍诱导有限平展满射 \(U'\to U\)。因此每个点 \(s\in S\) 都有既开又闭邻域 \(U\)，使 \(H_U\) 和 \(K_U\) 等平凡。断言 \SGARefLink{sga3:ix:proposition:2:11:item:ii:01}{(ii)} 立即得到。

\SGARefLink{sga3:ix:proposition:2:11:item:iii:01}{(iii)} 由 \SGARefLink{sga3:ix:proposition:2:11:item:i:01}{(i)}、\SGARefLink{sga3:ix:proposition:2:11:item:ii:01}{(ii)} 及定义立即得到。此外，拟等平凡的情形也将由 \SGARefLink{sga3:ix:comments:1:2}{1.2} 中预告的更一般事实“有限型 \(\Longrightarrow\) 拟等平凡”推出（参见 \SGARefLink{sga3:x:corollary:4-5}{X~4.5}）。

\SGARefTarget{sga3:ix:section:3}\section*{3. 无穷小性质：提升与共轭定理}
\addcontentsline{toc}{section}{3. 无穷小性质：提升与共轭定理}

乘法型群的基本无穷小性质由\SGARefLink{sga3:ix:theorem:3:1}{以下定理}得到。

\medskip
\noindent\SGARefTarget{sga3:ix:theorem:3:1}\textbf{定理 3.1.}
\textit{设 \(S\) 是仿射概形，\(H\) 是乘法型 \(S\)-群，\(\mathscr F\) 是 \(S\) 上赋有 \(H\) 作用的拟凝聚层（\SGARefLink{sga3:i:subsection:4-7}{I~4.7}）。则}
\[
  \mathrm H^i(H,\mathscr F)=0\qquad\text{对 }i>0,
\]
\textit{其中 \(\mathrm H^i\) 表示报告 I，第 5 节中研究的 Hochschild 上同调。}

% 源页：Exp9 组件本地第 8 页；完整阅读本第 654 页；原版印刷第 47--48 页。
事实上，由同上 \SGARefLink{sga3:ix:corollary:5:3}{5.3}，利用 \(S\) 和 \(G\) 的仿射环 \(A\) 和 \(B\)，以及定义 \(\mathscr F\) 的 \(A\)-模 \(M\)，可将 Hochschild 上同调表示为 \(A\)-模复形 \(C^\bullet(H,M)\) 的上同调；该复形的形成与每个基变换 \(A\to A'\) 交换。因此，对使 \(A'\) 在 \(A\) 上平坦的基变换 \(A\to A'\)，有
\[
  \mathrm H^i(H',\mathscr F')
  =\mathrm H^i(H,\mathscr F)\otimes_AA'.
\]
若 \(A'\) 还在 \(A\) 上忠实平坦，则为证明 \(\mathrm H^i(H,\mathscr F)\) 为零，只须证明 \(\mathrm H^i(H',\mathscr F')\) 为零。这将问题归结为验证 \(G\) 可对角化情形的 \SGARefLink{sga3:ix:theorem:3:1}{3.1}；该情形已在 \SGARefLink{sga3:i:theorem:5-3-3}{I~5.3.3} 中证明。

利用报告 III 的结果，我们将由 \SGARefLink{sga3:ix:theorem:3:1}{3.1} 推出若干几何方面的结论。

\medskip
\noindent\SGARefTarget{sga3:ix:theorem:3:2}\textbf{定理 3.2.}
\textit{设 \(S\) 是仿射概形，\(S_0\) 是由理想 \(J\) 定义的仿射子概形，\(H\) 是 \(S\) 上的乘法型群，\(G\) 是 \(S\) 上的任意群预概形。设 \(u,v:H\to G\) 是两个群同态，\(u_0,v_0:H_0\to G_0\) 是由 \(u,v\) 经基变换 \(S_0\to S\) 得到的同态。若 \(J^2=0\) 且 \(u_0=v_0\)，则存在 \(g\in G(S)\)，使}
\[
  v=\operatorname{int}(g)\circ u
  \qquad\textit{且}\qquad
  g_0=e.
\]

这由 \SGARefLink{sga3:iii:theorem:2-1:item:ii}{III~2.1(ii)} 和 \SGARefLink{sga3:ix:theorem:3:1}{3.1}（\(\mathrm H^1\) 的消没）推出。

\medskip
\noindent\SGARefTarget{sga3:ix:corollary:3:3}\textbf{推论 3.3.}
\textit{在 \SGARefLink{sga3:ix:theorem:3:2}{3.2} 的初始假设下，再假设 \(G\) 在 \(S\) 上平滑，但只要求 \(J\) 幂零（不必平方为零）。若存在 \(g_0\in G_0(S_0)\) 使 \(v_0=\operatorname{int}(g_0)\circ u_0\)，则 \(g_0\) 可提升为元素 \(g\in G(S)\)，使 \(v=\operatorname{int}(g)\circ u\)。}

对满足 \(J^n=0\) 的整数 \(n\) 作归纳，可归结到 \(J\) 平方为零的情形。由于 \(G\) 在 \(S\) 上平滑，可将 \(g_0\) 提升为 \(G(S)\) 的元素 \(g'\)。以 \(v'=\operatorname{int}(g')\circ u\) 代替 \(v\)，就归结到 \SGARefLink{sga3:ix:theorem:3:2}{3.2} 的情形。

\medskip
\noindent\SGARefTarget{sga3:ix:corollary:3:4}\textbf{推论 3.4.}
\textit{在 \SGARefLink{sga3:ix:theorem:3:2}{3.2} 的初始假设下，令 \(J\) 幂零，并进一步假设 \(u\) 是中心的（例如，\(G\) 交换）。则 \(u_0=v_0\) 蕴涵 \(u=v\)。}

事实上，同样可立即归结到 \(J\) 平方为零的情形，此时只须应用 \SGARefLink{sga3:ix:theorem:3:2}{3.2}。特别地：

\medskip
\noindent\SGARefTarget{sga3:ix:corollary:3:5}\textbf{推论 3.5.}
\textit{设 \(S,S_0,H,G\) 如 \SGARefLink{sga3:ix:theorem:3:2}{3.2}，且 \(J\) 幂零。设 \(u:H\to G\) 是 \(S\)-群同态，使 \(u_0:H_0\to G_0\) 为单位同态。则 \(u\) 为单位同态。}

\medskip
\noindent\SGARefTarget{sga3:ix:theorem:3:6}\textbf{定理 3.6.}
\textit{设 \(S\) 是仿射概形，\(S_0\) 是由幂零理想 \(J\) 定义的子概形，\(H\) 是 \(S\) 上的乘法型群，\(G\) 是 \(S\) 上平滑的群预概形，\(u_0:H_0\to G_0\) 是由基变换 \(S_0\to S\) 得到的 \(S_0\)-群的同态。则存在提升 \(u_0\) 的 \(S\)-群同态 \(u:H\to G\)。此外，由 \SGARefLink{sga3:ix:corollary:3:3}{3.3}，任何两个这样的提升 \(u,u'\) 都由 \(G(S)\) 中某个在 \(G(S_0)\) 中约化为单位元的元素共轭。}

% 源页：Exp9 组件本地第 9 页；完整阅读本第 655 页；原版印刷第 49 页。
这由 \SGARefLink{sga3:iii:theorem:2-1}{III~2.1}、\SGARefLink{sga3:ix:proposition:2:3}{2.3} 和 \SGARefLink{sga3:ix:theorem:3:1}{3.1} 推出：\(\mathrm H^2\) 的消没给出 \(u_0\) 的提升的存在性，\(\mathrm H^1\) 的消没给出相差共轭意义下的唯一性。

可用相同方法证明 \SGARefLink{sga3:ix:theorem:3:2}{3.2}--\SGARefLink{sga3:ix:theorem:3:6}{3.6} 的以下变式。事实上，将这些变式应用于那里所考虑同态的图子群，即可推出前述陈述。

\medskip
\noindent\SGARefTarget{sga3:ix:theorem:3:2:bis}\textbf{定理 3.2 bis.}
\textit{设 \(S\) 是仿射概形，\(S_0\) 是由理想 \(J\) 定义的子概形，\(G\) 是 \(S\)-群预概形，\(H,H'\) 是乘法型子群预概形。设 \(G_0,H_0,H'_0\) 为经基变换 \(S_0\to S\) 得到的 \(S_0\) 上的群。若 \(J^2=0\) 且 \(H_0=H'_0\)，则存在 \(g\in G(S)\)，使}
\[
  H'=\operatorname{int}(g)(H)
  \qquad\textit{且}\qquad
  g_0=e.
\]

\medskip
\noindent\SGARefTarget{sga3:ix:corollary:3:3:bis}\textbf{推论 3.3 bis.}
\textit{在 \SGARefLink{sga3:ix:theorem:3:2:bis}{3.2 bis} 的初始假设下，再假设 \(G\) 在 \(S\) 上平滑，且 \(J\) 幂零（不必平方为零）。若 \(g_0\in G_0(S_0)\) 满足 \(H'_0=\operatorname{int}(g_0)(H_0)\)，则 \(g_0\) 可提升为元素 \(g\in G(S)\)，使 \(H'=\operatorname{int}(g)(H)\)。}

\medskip
\noindent\SGARefTarget{sga3:ix:corollary:3:4:bis}\textbf{推论 3.4 bis.}
\textit{在 \SGARefLink{sga3:ix:theorem:3:2:bis}{3.2 bis} 的初始假设下，假设 \(J\) 幂零，且 \(H\) 是中心的（例如，\(G\) 交换）。则 \(H_0=H'_0\) 蕴涵 \(H=H'\)。}

\medskip
\noindent\SGARefTarget{sga3:ix:theorem:3:6:bis}\textbf{定理 3.6 bis.}
\textit{设 \(S\) 是仿射概形，\(S_0\) 是由幂零理想 \(J\) 定义的子概形，\(G\) 是在 \(S\) 上平滑的 \(S\)-群预概形，\(H_0\) 是 \(G_0=G\times_SS_0\) 的乘法型子群预概形。则：}
\begin{enumerate}
  \item[(a)]\SGARefTarget{sga3:ix:theorem:3:6:bis:item:a:01}
  \textit{存在 \(G\) 的子群预概形 \(H\)，在 \(S\) 上平坦，且 \(H\times_SS_0=H_0\)；}

  \item[\SGAInlineRefTarget{sga3:ix:theorem:3:2:item:b}(b)]\SGARefTarget{sga3:ix:theorem:3:6:bis:item:b:01}
  \textit{这样的 \(H\) 必为乘法型；}

  \item[(c)]\SGARefTarget{sga3:ix:theorem:3:6:bis:item:c:01}
  \textit{最后，由 \SGARefLink{sga3:ix:theorem:3:2:bis}{3.2 bis}，\(H_0\) 的任何两个这样的提升 \(H,H'\) 都由某个在 \(G(S_0)\) 中约化为单位元的元素 \(g\in G(S)\) 共轭。}
\end{enumerate}

下面说明如何证明 \SGARefLink{sga3:ix:theorem:3:2:bis}{3.2 bis}（\SGARefLink{sga3:ix:corollary:3:3:bis}{3.3 bis} 和 \SGARefLink{sga3:ix:corollary:3:4:bis}{3.4 bis} 是它的直接推论）以及 \SGARefLink{sga3:ix:theorem:3:6:bis}{3.6 bis} 的断言 \SGARefLink{sga3:ix:theorem:3:6:bis:item:a:01}{(a)}。后者由 \SGARefLink{sga3:ix:theorem:3:1}{3.1}（\(\mathrm H^2\) 的消没）及 \SGARefLink{sga3:iii:corollary:4-37:item:ii}{III~4.37(ii)} 得到。类似地，\SGARefLink{sga3:ix:theorem:3:2:bis}{3.2 bis} 由 \SGARefLink{sga3:ix:theorem:3:1}{3.1}（\(\mathrm H^1\) 的消没）及 \SGARefLink{sga3:iii:corollary:4-37:item:i}{III~4.37(i)} 得到，至少在 \(G\) 和 \(H\) 均在 \(S\) 上平滑时如此。

不过，利用下一篇报告的结果，可非常简单地由 \SGARefLink{sga3:ix:theorem:3:2}{3.2}、\SGARefLink{sga3:ix:theorem:3:6}{3.6} 分别推出此处所述一般形式的 \SGARefLink{sga3:ix:theorem:3:2:bis}{3.2 bis}、\SGARefLink{sga3:ix:theorem:3:6:bis}{3.6 bis}。对于 \SGARefLink{sga3:ix:theorem:3:2:bis}{3.2 bis}，只须注意，由 \SGARefLink{sga3:x:proposition:2:1}{X~2.1}，\(H\) 与 \(H'\) 同构，于是归结到 \SGARefLink{sga3:ix:theorem:3:2}{3.2}。

对于 \SGARefLink{sga3:ix:theorem:3:6:bis}{3.6 bis}，注意 \(H_0\) 必在 \(S_0\) 上有限型：\footnote{编者注：原文说，“由 \SGARefLink{sga3:ix:proposition:2:1:item:b:01}{2.1(b)}，因为其纤维有限型”。编者未能理解此论证，将其替换为以下论证。}事实上，作为在 \(S_0\) 上平滑的 \(G_0\) 的子群预概形，群 \(H_0\) 在 \(S_0\) 上局部有限型。但由 \SGARefLink{sga3:ix:proposition:2:1:item:a:01}{2.1(a)}，\(H_0\) 在 \(S_0\) 上仿射，因而在 \(S_0\) 上有限型。所以由 \SGARefLink{sga3:x:corollary:4-5}{X~4.5}，\(H_0\) 拟等平凡，从而由 \SGARefLink{sga3:x:proposition:2:1}{X~2.1}，它来自 \(S\) 上的乘法型群 \(H\)。\footnote{编者注：下句已扩充。}由 \SGARefLink{sga3:ix:theorem:3:6}{3.6}，存在提升浸入 \(H_0\hookrightarrow
G_0\) 的 \(S\)-群同态 \(u:H\to G\)。由于 \(H\) 和 \(H_0\)（分别地，\(G\) 和 \(G_0\)）具有相同的底层拓扑空间，且对每个 \(h\in H\)，态射

% 源页：Exp9 组件本地第 10 页；完整阅读本第 656 页；原版印刷第 50 页。
\[
  \mathscr O_{G,u(h)}\longrightarrow\mathscr O_{H,h}
\]
是满射（因为模幂零理想 \(J\) 约化后它是满射），故态射 \(u\) 也是浸入（参见 EGA~I，4.2.2）。

最后，由 \SGARefLink{sga3:x:corollary:2:3}{X~2.3}，\(H_0\) 到 \(G\) 的平坦子群的每个提升 \(H\) 必为乘法型，\footnote{编者注：\SGARefLink{sga3:x:corollary:2:3}{X~2.3} 本质上依赖本报告的 \SGARefLink{sga3:ix:theorem:3:6}{定理~3.6}。}这也证明了 \SGARefLink{sga3:ix:theorem:3:6:bis}{3.6 bis} 的断言 \SGARefLink{sga3:ix:theorem:3:6:bis:item:b:01}{(b)}。读者可验证，报告 X 并未使用结果 \SGARefLink{sga3:ix:theorem:3:2:bis}{3.2 bis}--\SGARefLink{sga3:ix:theorem:3:6:bis}{3.6 bis}，因此不存在循环论证。

\medskip
\noindent\SGARefTarget{sga3:ix:proposition:3:8}\textbf{命题 3.8.}
\footnote{编者注：保留原文编号；没有 3.7。}
\textit{设 \(S\) 是预概形，\(S_0\) 是由局部幂零理想 \(J\) 定义的闭子概形，\(G\) 是乘法型 \(S\)-群，\(X\) 是局部有限呈示的 \(S\)-预概形。假设 \(G\) 作用于 \(X\)，使 \(G_0=G\times_SS_0\) 在 \(X_0=X\times_SS_0\) 上的作用平凡。则 \(G\) 在 \(X\) 上的作用平凡。}

其证明就是 \SGARefLink{sga3:iii:corollary:2-4:item:b}{III~2.4(b)} 的证明。也可注意到 \SGARefLink{sga3:ix:proposition:3:8}{3.8} 立即归结到 \(G\) 可对角化的情形，而该处处理的正是这种情形，从而归结到所引结果。

\medskip
\noindent\SGARefTarget{sga3:ix:corollary:3:9}\textbf{推论 3.9.}
\textit{设 \(S,S_0\) 如上，\(u:G\to H\) 是 \(S\)-群同态，其中 \(G\) 是乘法型，\(H\) 在 \(S\) 上局部有限呈示。假设经基变换 \(S_0\to S\) 得到的同态 \(u_0:G_0\to H_0\) 是中心的。则 \(u\) 是中心的。}

事实上，只须取 \(X=H\) 应用 \SGARefLink{sga3:ix:proposition:3:8}{3.8}，令 \(G\) 按下式作用于 \(H\)：
\[
  (g,h)\longmapsto\operatorname{int}(u(g))\cdot h
  =u(g)h\,u(g)^{-1}.
\]

% 终止位置：第 4 节之前，Exp9 组件本地第 10 页
% （完整阅读本第 656 页；原版印刷第 50 页）。
